% sec-02: the ten recipients, their mechanisms, and their asks.
% Resolved from applications/README.md and the ten app-*/README.md files.
% Figures 3 (d2 grid) and 6 (graphviz DAG).
\section{The Ten Applications}

\subsection{Selection Rule}

A recipient qualifies for this set only if \textit{Science: A New Golden Age}
either names its programme or names the mechanism that programme runs. Nothing
was included on plausibility. The rule matters because it is what makes the set
an argument rather than a mailing list: every application can open by quoting the
recipient's own mechanism back to it, in the administration's wording, and then
answer it.

Three mechanism families emerged from applying that rule. \textbf{Person-based},
where the report names the NIH Director's Pioneer Award by title and cites
HHMI-style support as the evidence. \textbf{Organization-type}, where the report
names X-Labs as the first federal programme designed to fund independent
research organizations outside traditional academia, names focused research
organizations as time-bound nonprofits built to dissolve, names the ARPA
program-manager model, and names SBIR as opening doors for technician-founded
ventures. \textbf{Partnership and cost share}, where the report holds up the
Foundation for the NIH and the Accelerating Medicines Partnership as its model,
names the Genesis Mission and its Robotics line, and asks agencies to prioritise
non-federal cost share \cite{kratsiosvought2026fy2028priorities}.

\begin{appfloat}[!tb]
\begin{appfig}[d2-type / scored grid]
% Five columns on a fixed pitch, placed by anchor from the left neighbour so
% they cannot drift; eleven rows of equal minimum height.  Column widths are
% 4.0, 2.9, 2.5, 2.2 and 2.3 cm, totalling 13.9 cm.
\node[d2cellh,text width=3.6cm,minimum width=4.0cm,minimum height=6.0mm] (h0) at (0,0) {Recipient};
\node[d2cellh,text width=2.5cm,minimum width=2.9cm,minimum height=6.0mm,anchor=west] (h1) at (h0.east) {Mechanism};
\node[d2cellh,text width=2.1cm,minimum width=2.5cm,minimum height=6.0mm,anchor=west] (h2) at (h1.east) {Term};
\node[d2cellh,text width=1.8cm,minimum width=2.2cm,minimum height=6.0mm,anchor=west] (h3) at (h2.east) {Direct ask};
\node[d2cellh,text width=1.9cm,minimum width=2.3cm,minimum height=6.0mm,anchor=west] (h4) at (h3.east) {Lead};
\foreach \r/\prev/\rec/\mech/\term/\ask/\lead/\lf in {%
  1/h0/{01 NIH Pioneer Award}/{person-based}/{5 years}/{\$3.5M}/{surgical}/d2soft,
  2/r1/{02 ARPA-H}/{milestone}/{36 months}/{\$2.1M}/{surgical}/d2soft,
  3/r2/{03 NSF TIP X-Labs}/{organization}/{5 years}/{\$3.5M}/{surgical}/d2soft,
  4/r3/{04 DOE Genesis Mission}/{mission}/{5 years}/{\$3.5M}/{surgical}/d2soft,
  5/r4/{05 NIH SEED SBIR}/{small business}/{9 + 24 mo}/{\$1.6M}/{surgical}/d2soft,
  6/r5/{06 FNIH AMP}/{consortium}/{5 years}/{\$3.5M}/{medical}/d2cellg,
  7/r6/{07 HHMI Investigator}/{person-based}/{7 years}/{\$0.7M/yr}/{medical}/d2cellg,
  8/r7/{08 NCI CTEP}/{concept review}/{5 years}/{\$3.5M}/{medical}/d2cellg,
  9/r8/{09 Convergent FRO}/{time-bound}/{5 years}/{\$3.5M}/{medical}/d2cellg,
  10/r9/{10 UC San Diego Moores}/{institutional}/{one meeting}/{none}/{medical}/d2cellg}{%
  \node[d2celll,text width=3.6cm,minimum width=4.0cm,minimum height=6.0mm,anchor=north west]
    (r\r) at (\prev.south west) {\rec};
  \node[d2cell,text width=2.5cm,minimum width=2.9cm,minimum height=6.0mm,anchor=west]
    (m\r) at (r\r.east) {\mech};
  \node[d2cell,text width=2.1cm,minimum width=2.5cm,minimum height=6.0mm,anchor=west]
    (t\r) at (m\r.east) {\term};
  \node[d2cell,text width=1.8cm,minimum width=2.2cm,minimum height=6.0mm,anchor=west]
    (a\r) at (t\r.east) {\ask};
  \node[\lf,text width=1.9cm,minimum width=2.3cm,minimum height=6.0mm,anchor=west]
    (v\r) at (a\r.east) {\lead};}
\node[font=\tiny\itshape,text=protogray,anchor=north west,text width=136mm]
  at ([yshift=-2.5mm]r10.south west) {The science column is omitted because it
  would be constant: the same population, the same 3+3 escalation, the same
  co-primary endpoints, and the same advisory boundary appear in all ten.};
\end{appfig}
\vspace{-0.7cm}
\figcaption{Ten applications on four attributes. The science column is constant
by construction: the same trial is described ten times, and only the mechanism,
the term, and the ask change.}
\end{appfloat}

\subsection{What Differs, and What Does Not}

\begin{apptable}
\begin{tabularx}{\textwidth}{>{\raggedright\arraybackslash}p{0.7cm} >{\raggedright\arraybackslash}p{3.4cm} >{\raggedright\arraybackslash}p{2.6cm} >{\raggedright\arraybackslash}X}
\headrow \thead{No} & \thead{Recipient} & \thead{Mechanism family} & \thead{The anchor sentence in the report}\\
\midrule
01 & NIH Common Fund, Director's Pioneer Award & Person-based & ``scaling long-horizon grants for the best and brightest modeled on \dots\ Director's Pioneer Award'' \\
02 & ARPA-H mission office & Organization-type & The ARPA program-manager model, with ARPA-H and ARPA-E as the live examples \\
03 & NSF TIP Directorate, X-Labs & Organization-type & ``the first federal program designed to fund independent research organizations outside traditional academia'' \\
04 & DOE Office of Science, Genesis Mission & Partnership & Robotics mission: ``to initiate the era of physical AI-driven scientific discovery'' \\
05 & NIH SEED, SBIR and STTR & Organization-type & ``Programs like SBIR open doors for technician-founded ventures'' \\
06 & Foundation for the NIH, AMP & Partnership & FNIH and the Accelerating Medicines Partnership held up as the model \\
07 & HHMI Investigator Program & Person-based & Person-based funding of roughly \$10 million over seven years \\
08 & NCI Cancer Therapy Evaluation Program & Partnership & Prioritised foundational biological sciences; the report's own cancer framing \\
09 & Convergent Research, FRO programme & Organization-type & Focused research organizations as time-bound nonprofit research startups \\
10 & UC San Diego Moores Cancer Center & Partnership & Regional innovation clusters; non-federal cost share \\
\end{tabularx}
\end{apptable}
\tabcap{The ten recipients with the sentence in the report that qualifies each.
No recipient appears without an anchor, which is what makes the set an argument
inside the funder's own frame rather than a mailing list.}

What is constant is the trial. Every one of the ten describes the same
population, resectable or borderline-resectable KRAS G12-mutated pancreatic
ductal adenocarcinoma; the same design, open-label, single-arm, 3+3 escalation,
up to 18 treated participants; the same co-primary endpoints; and the same
advisory boundary, with arm-level stop specified at 3 milliseconds and
system-wide stop at 500 milliseconds \cite{onpremwhippl,phase1ind}.

That constancy is the point. A reviewer who compares two applications from this
set is comparing two funding mechanisms applied to one piece of science, with
the science held fixed. It is the closest thing an individual applicant can
construct to a controlled comparison of mechanisms, and it is only possible
because the protocol was published before the applications were written.

\subsection{Redundancy Across the Set}

Sending ten applications is not ten times one application. Because the ten fund
overlapping activity sets, most of the programme survives a refusal from any one
recipient. Figure 6 makes that concrete, and identifies the two activities that
do not.

\begin{appfloat}[!tb]
\begin{appfig}[graphviz-type / dependency DAG]
% Ten awards in two columns of five at x = 0 and 2.4, pitch 1.15, so the left
% cluster matches the right in height and the edge fan does not cross itself.
% Eight activities at x = 8.6, pitch 1.30.  The corridor between x = 4.2 and
% x = 7.4 carries every edge and holds no node.
\foreach \i/\x/\y/\lab in {1/0/2.3/{01 Pioneer}, 2/0/1.15/{02 ARPA-H}, 3/0/0/{03 X-Labs},
  4/0/-1.15/{04 Genesis}, 5/0/-2.3/{05 SBIR}, 6/2.5/2.3/{06 FNIH}, 7/2.5/1.15/{07 HHMI},
  8/2.5/0/{08 CTEP}, 9/2.5/-1.15/{09 FRO}, 10/2.5/-2.3/{10 Moores}}{
  \node[gvnode,text width=17mm] (w\i) at (\x,\y) {\lab};}
\begin{scope}[on background layer]
\node[gvcluster,fit=(w1)(w5)(w6)(w10),inner sep=7pt] (cl) {};
\end{scope}
\node[gvctitle,anchor=south west] at ([yshift=1mm]cl.north west) {Ten applications};

\foreach \i/\y/\lab/\sty in {1/3.0/{Site agreement}/gvbox, 2/1.7/{IND activation}/gvboxk,
  3/0.4/{IRB approval}/gvboxk, 4/-0.9/{Interlock rig}/gvbox, 5/-2.2/{Logging schema}/gvbox,
  6/-3.5/{Cohort accrual}/gvbox}{
  \node[\sty,text width=24mm] (t\i) at (9.0,\y) {\lab};}
\node[gvbox,text width=24mm] (t7) at (9.0,-4.8) {Specimen pipeline};
\node[gvbox,text width=24mm] (t8) at (9.0,-6.1) {Public release};
\begin{scope}[on background layer]
\node[gvcluster2,fit=(t1)(t8),inner sep=7pt] (ca) {};
\end{scope}
\node[gvctitle,anchor=south west] at ([yshift=1mm]ca.north west) {Trial activities};

\foreach \s/\t in {w1/t1, w1/t2, w1/t6, w2/t2, w2/t4, w2/t6, w3/t4, w3/t5, w3/t8,
  w4/t5, w4/t8, w5/t4, w5/t5, w6/t7, w6/t6, w7/t6, w7/t7, w8/t6, w8/t7, w8/t8,
  w9/t5, w9/t7, w9/t8, w10/t1, w10/t3}{
  \draw[gvedge] (\s.east) -- (\t.west);}
\node[font=\tiny\itshape,text=protogray,anchor=north,text width=118mm,align=center]
  at (5.6,-7.2) {IND activation has two upstream awards and IRB approval has
  one. Both are drawn heavier, and both are regulatory rather than technical,
  which is why every application that can fund them funds them in year one.};
\end{appfig}
\vspace{-0.7cm}
\figcaption{Ten awards against eight trial activities. Six activities have two or
more upstream awards and survive any single refusal; the two that do not are
drawn heavier, and both are regulatory.}
\end{appfloat}

Six of the eight activities have two or more upstream awards. The interlock rig
is funded by applications 02, 03, and 05; the logging schema by 03, 04, 05, and
09; cohort accrual by 01, 02, 06, 07, and 08. Any single refusal in that set
leaves the activity funded.

The two exceptions are regulatory. IND activation is reachable only through
applications 01 and 02, and IRB approval only through application 10, which asks
for no money at all. That asymmetry is the programme's real structure: the
scarce resource is not funding but institutional standing, and the one
application that could supply it is the one with nothing to offer in return.
